\section{All Nodes at Distance K from a Target}
\label{sec:trees-distance-k}

\problemstatement{Given the root, a target node, and an integer $K$, return
the values of all nodes that are exactly $K$ edges away from the target.
Distance may go upward (toward ancestors) as well as downward.}

\begin{patternbox}
The twist is \emph{``distance can go upward.''} A binary tree node has no
parent pointer, so radiating outward in all directions is impossible with
the raw structure. The fix: \textbf{build a parent map} so every node
effectively becomes a graph vertex with up-to-three neighbours (left
child, right child, parent), then run a plain \textbf{BFS} from the target.
\end{patternbox}

\subsection*{Understanding the problem carefully}

``All nodes at distance $K$'' is a BFS problem -- but BFS needs to move in
every direction, including up to the parent. So we first walk the tree
once recording each node's parent in a hash map. Now the tree behaves like
an undirected graph: from any node we can step to its left child, right
child, or parent. A breadth-first expansion from the target, stopping after
$K$ layers, lands exactly on the distance-$K$ nodes.

\begin{ideabox}[Turn the Tree into a Graph, Then BFS $K$ Layers]
Record parents with one DFS. Start BFS from the target with a
\textit{visited} set (so we never walk back the way we came). Expand
layer by layer; after $K$ layers the queue holds precisely the nodes at
distance $K$.
\end{ideabox}

\subsection*{Algorithm}

\begin{enumerate}
  \item DFS to fill \textit{parent}[node] for every node.
  \item BFS from target: mark it visited, distance $0$.
  \item Repeat $K$ times: expand the current frontier to all unvisited
        neighbours (left, right, parent), marking them visited.
  \item After $K$ expansions, the frontier is the answer.
\end{enumerate}

\subsection*{Dry run}

\dryrun{Tree: root $3$; left $5$ (children $6,2$); right $1$. Target $5$,
$K=2$.}

\begin{itemize}
  \item Parents: $5\!\to\!3$, $6\!\to\!5$, $2\!\to\!5$, $1\!\to\!3$.
  \item Layer $0$: $\{5\}$. Layer $1$: neighbours of $5$ = $6,2,3$
        (children and parent).
  \item Layer $2$: from $6$ nothing new; from $2$ nothing new; from $3$
        neighbours are $5$ (visited) and $1$ (new). Frontier $=\{1\}$.
  \item Nodes at distance $2$: $1$ (and, in the classic example, $7,4$ if
        present under $1$).
\end{itemize}

\subsection*{C++ implementation}

\begin{lstlisting}
#include <bits/stdc++.h>
using namespace std;

struct TreeNode {
    int val;
    TreeNode* left;
    TreeNode* right;
    TreeNode(int v) : val(v), left(nullptr), right(nullptr) {}
};

void markParents(TreeNode* node, unordered_map<TreeNode*, TreeNode*>& par) {
    if (node == nullptr) return;
    if (node->left)  { par[node->left]  = node; markParents(node->left,  par); }
    if (node->right) { par[node->right] = node; markParents(node->right, par); }
}

vector<int> distanceK(TreeNode* root, TreeNode* target, int K) {
    unordered_map<TreeNode*, TreeNode*> par;
    markParents(root, par);

    unordered_set<TreeNode*> visited;
    queue<TreeNode*> q;
    q.push(target);
    visited.insert(target);

    int dist = 0;
    while (!q.empty()) {
        if (dist == K) break;
        int s = q.size();
        for (int i = 0; i < s; i++) {
            TreeNode* node = q.front(); q.pop();
            for (TreeNode* nb : {node->left, node->right, par.count(node) ? par[node] : nullptr}) {
                if (nb && !visited.count(nb)) { visited.insert(nb); q.push(nb); }
            }
        }
        dist++;
    }

    vector<int> ans;
    while (!q.empty()) { ans.push_back(q.front()->val); q.pop(); }
    return ans;
}

int main() {
    TreeNode* root = new TreeNode(3);
    root->left = new TreeNode(5);
    root->right = new TreeNode(1);
    root->left->left  = new TreeNode(6);
    root->left->right = new TreeNode(2);

    vector<int> res = distanceK(root, root->left, 2);
    sort(res.begin(), res.end());
    for (int x : res) cout << x << " ";
    cout << "\n";
    return 0;
}
\end{lstlisting}

\begin{complexitybox}
\textbf{Time Complexity.} $O(n)$ -- one DFS to record parents and a BFS
visiting each node at most once. \textbf{Space Complexity.} $O(n)$ for the
parent map, visited set, and queue.
\end{complexitybox}

\begin{takeawaybox}
``Distance in all directions'' on a tree is a graph-BFS problem in
disguise; the enabling move is manufacturing parent pointers so upward
travel becomes possible. This parent-map trick recurs in the next section's
burning-tree problem.
\end{takeawaybox}
